\documentclass[11pt, a4paper]{article}

% Packages
\usepackage[margin=2.5cm]{geometry}
\usepackage{booktabs}
\usepackage{array}
\usepackage{longtable}
\usepackage{listings}
\lstset{
	basicstyle=\scriptsize\ttfamily,
}
\usepackage{xcolor}
\usepackage{hyperref}
\usepackage{parskip}
\usepackage{enumitem}
\usepackage{titlesec}
\usepackage{microtype}
\usepackage{amsmath}   % for \text{} in math mode
\usepackage{amssymb}  % for checkmark

% Hyperlink styling
\hypersetup{
  colorlinks = true,
  linkcolor  = blue!60!black,
  urlcolor   = blue!60!black,
}

% Code listing style (for R code blocks)
\lstdefinestyle{Rstyle}{
  language        = R,
  basicstyle      = \ttfamily\small,
  keywordstyle    = \color{blue!70!black}\bfseries,
  commentstyle    = \color{green!50!black}\itshape,
  stringstyle     = \color{orange!80!black},
  backgroundcolor = \color{gray!8},
  frame           = single,
  framerule       = 0.4pt,
  rulecolor       = \color{gray!40},
  breaklines      = true,
  breakatwhitespace = true,
  showstringspaces = false,
  tabsize         = 2,
  captionpos      = b,
  xleftmargin     = 6pt,
  xrightmargin    = 6pt,
  aboveskip       = 8pt,
  belowskip       = 8pt,
}

\lstset{style=Rstyle}

% Section formatting
\titleformat{\section}{\large\bfseries}{}{0em}{}[\titlerule]
\titleformat{\subsection}{\normalsize\bfseries}{}{0em}{}
\titleformat{\subsubsection}{\normalsize\itshape\bfseries}{}{0em}{}

% Checkmark command
\newcommand{\done}{$\checkmark$}

% -----------------------------------------------------------------------
\begin{document}

\scriptsize

\begin{center}
  {\Large\bfseries Project Memory Log: Morphological Processing ERP Paper}\\[4pt]
  {\normalsize Conversation Summary and Continuity Document}
\end{center}

\medskip
\hrule
\bigskip

% -----------------------------------------------------------------------
\section{1.\ The Original Problem}

The user is writing a specialist psycholinguistics journal paper examining individual differences in morphological processing using ERPs and reaction times. The study involves a lexical decision task (LDT) with:

\begin{itemize}[leftmargin=1.5em]
  \item \textbf{Participants:} 67 adults (60 with ERP data)
  \item \textbf{Stimuli:} 100 words and 100 morphologically complex nonwords
  \item \textbf{Nonword types:} Complex (real stem + real affix, e.g., \textit{gasful}) vs.\ Simple (real stem + pseudoaffix, e.g., \textit{gasfil})
  \item \textbf{Measures:} Reaction times (RTs) and ERPs (N250 and N400 components)
  \item \textbf{Individual differences:} Two orthogonal PCA dimensions derived from vocabulary knowledge, spelling ability, and print exposure:
  \begin{itemize}
    \item \textbf{Dim.1 = Language Proficiency (LP)} — depth/breadth of lexical-semantic knowledge
    \item \textbf{Dim.2 = Orthographic Sensitivity (OS)} — sensitivity to sublexical orthographic structure
  \end{itemize}
\end{itemize}

The user was struggling with: (a) theoretical framing of the introduction, (b) analytic strategy, (c) R code for mixed effects models, and (d) writing up results.

% -----------------------------------------------------------------------
\section{2.\ Key Theoretical Decisions}

\subsection{The Core Prediction: Double Dissociation}

The study tests whether OS and LP dissociably modulate morphological processing at distinct temporal stages:
\begin{itemize}[leftmargin=1.5em]
  \item \textbf{OS $\to$ early form-based processing (N250):} OS should modulate stem/base frequency effects at the N250, reflecting sensitivity to sublexical orthographic structure.
  \item \textbf{LP $\to$ later semantic integration (N400):} LP should modulate family size effects at the N400, reflecting depth of lexical-semantic knowledge.
\end{itemize}

This prediction follows from Andrews and colleagues' finding that orthographic and semantic profiles are dissociable axes of individual variation, combined with the theoretical distinction between stem frequency (form-level access) and family size (semantic representation richness).

\subsection{The Three-Question Analytic Structure}

The analyses are organised around three core questions:
\begin{enumerate}[leftmargin=1.5em]
  \item \textbf{Q1 (RT/Behavioral):} Does OS influence overall lexical decision efficiency?
  \item \textbf{Q2 (N250/Early neural):} Does OS modulate early morpho-orthographic parsing?
  \item \textbf{Q3 (N400/Late neural):} Does LP shape semantic integration of morphological information?
\end{enumerate}

\subsection{Key Theoretical Models Discussed}

\begin{itemize}[leftmargin=1.5em]
  \item \textbf{Baayen--Schreuder race model} (classical dual-route): stem frequency $\to$ decomposition route; family size $\to$ whole-form route
  \item \textbf{Early automatic decomposition (Rastle \& Davis):} obligatory early form-based parse, then semantic filter
  \item \textbf{Morpho-orthographic account (Grainger):} intermediate morpho-orthographic level
  \item \textbf{NDL/LDL (Baayen et al.):} no morpheme representations; discriminative learning from distributional statistics; both effects emergent from learned weight structure
\end{itemize}

% -----------------------------------------------------------------------
\section{3.\ Introduction — Current Status: Complete}

\subsection{Structure}
\begin{enumerate}[leftmargin=1.5em]
  \item Opening paragraphs on reading and English morphology
  \item Bridge paragraph planting individual differences logic
  \item \textbf{Subsection: Stem Frequency and Morphological Family Size: Theoretical Interpretations Across Current Models} — covers classical, early decomposition, morpho-orthographic, and NDL/LDL accounts
  \item \textbf{Subsection: Individual Differences in Morphological Processing} — Lexical Quality Hypothesis; Andrews et al.\ (2013) finding; OS/LP profiles mapped onto theoretical accounts
  \item \textbf{Subsection: Aims and Predictions} — three analyses, double dissociation as central claim
\end{enumerate}

\subsection{Key Framing Decisions}

\begin{itemize}[leftmargin=1.5em]
  \item Words and nonwords framed as doing distinct theoretical work (words = time course of stored representations; nonwords = generalisation to novel forms)
  \item RT models framed as efficiency measures (no morphological $\times$ ID interactions by design — stated explicitly in methods)
  \item ERP models framed as architecture measures
  \item The asymmetric interaction structure (OS at N250, LP at N400) is justified in both the aims section and the results section
\end{itemize}

% -----------------------------------------------------------------------
\section{4.\ Analytic Strategy — Finalised}

\subsection{RT Models}

No morphological structure $\times$ individual difference interactions in RT models by design. Individual differences entered as main effects only.

\textbf{Words RT:}
\begin{lstlisting}
m_rt_words <- lmer(
  logRT ~ zLogBF + zLogFS + OS + LP +
    (1 | SubjID) + (1 | Item),
  data = rt_w,
  control = lmerControl(optimizer = "bobyqa")
)
\end{lstlisting}

\textbf{Nonwords RT:}
\begin{lstlisting}
m_rt_nw <- lmer(
  logRT ~ complexity + zLogBF + zLogFS + OS + LP +
    (1 + complexity | SubjID) + (1 | ItemID),
  data = rt_nw,
  control = lmerControl(optimizer = "bobyqa")
)
\end{lstlisting}

\subsection{Contrast Coding}

\begin{itemize}[leftmargin=1.5em]
  \item \texttt{complexity}: sum coded \texttt{c(-0.5, 0.5)} with levels \texttt{c("Simple", "Complex")}
  \item \texttt{family\_size}: sum coded \texttt{c(-0.5, 0.5)} with levels \texttt{c("Small", "Large")}
  \item \texttt{base\_frequency}: sum coded \texttt{c(-0.5, 0.5)} with levels \texttt{c("Low", "High")}
  \item OS and LP: PCA scores, already mean-centred by construction — no additional scaling needed
\end{itemize}

\subsection{ERP Models — Base Frequency Analyses (Primary for Words)}

\textbf{Words N250 (Base Frequency $\times$ OS):}
\begin{lstlisting}
contrasts(erp_w$base_frequency) <- c(-0.5, 0.5)
m_w_N250_bf <- lmer(
  value ~ base_frequency * OS + LP +
    (1 | SubjID) + (1 | SubjID:chlabel),
  data = subset(erp_w, time_window == "N250"),
  control = lmerControl(optimizer = "bobyqa")
)
\end{lstlisting}

\textbf{Words N400 (Base Frequency $\times$ LP):}
\begin{lstlisting}
m_w_N400_bf <- lmer(
  value ~ base_frequency * LP + OS +
    (1 | SubjID) + (1 | SubjID:chlabel),
  data = subset(erp_w, time_window == "N400"),
  control = lmerControl(optimizer = "bobyqa")
)
\end{lstlisting}

\subsection{ERP Models — Family Size Analyses}

\textbf{Words N250 (Family Size $\times$ OS) — secondary:}
\begin{lstlisting}
contrasts(erp_w$family_size) <- c(-0.5, 0.5)
m_w_N250_fs <- lmer(
  value ~ family_size * OS + LP +
    (1 | SubjID) + (1 | SubjID:chlabel),
  data = subset(erp_w, time_window == "N250"),
  control = lmerControl(optimizer = "bobyqa")
)
\end{lstlisting}

\textbf{Words N400 (Family Size $\times$ LP) — secondary:}
\begin{lstlisting}
m_w_N400_fs <- lmer(
  value ~ family_size * LP + OS +
    (1 | SubjID) + (1 | SubjID:chlabel),
  data = subset(erp_w, time_window == "N400"),
  control = lmerControl(optimizer = "bobyqa")
)
\end{lstlisting}

\textbf{Nonwords N250 (Complexity $\times$ OS and LP; \texttt{family\_size} as covariate — confirmed final):}
\begin{lstlisting}
contrasts(erp_nw$family_size) <- c(-0.5, 0.5)
m_nw_N250 <- lmer(
  value ~ complexity * OS + complexity * LP + family_size +
    (1 | SubjID) + (1 | SubjID:chlabel),
  data = subset(erp_nw, time_window == "N250"),
  control = lmerControl(optimizer = "bobyqa")
)
\end{lstlisting}

\textbf{Nonwords N400 (Complexity $\times$ Family Size $\times$ LP — centerpiece finding):}
\begin{lstlisting}
contrasts(erp_nw$complexity)  <- c(-0.5, 0.5)
contrasts(erp_nw$family_size) <- c(-0.5, 0.5)
m_nw_N400 <- lmer(
  value ~ complexity * family_size * LP + OS +
    (1 | SubjID) + (1 | SubjID:chlabel),
  data = subset(erp_nw, time_window == "N400"),
  control = lmerControl(optimizer = "bobyqa")
)
\end{lstlisting}

\subsection{R Packages Used}

\begin{lstlisting}
library(lme4)
library(lmerTest)    # Satterthwaite df -- load BEFORE lmerTest, BEFORE fitting models
library(emmeans)     # follow-up contrasts; use emm_options(lmer.df = "kenward-roger")
library(effectsize)  # eta_squared(model, partial = TRUE)
library(performance) # r2_nakagawa(model) -- use this, NOT r2()
library(MuMIn)       # r.squaredGLMM() as alternative if performance fails
library(ggplot2)
library(patchwork)   # load after ggplot2, before lmerTest
\end{lstlisting}

\subsection{Important Technical Notes}

\begin{itemize}[leftmargin=1.5em]
  \item Always load \texttt{ggplot2} and \texttt{patchwork} \textit{before} \texttt{lmerTest} to avoid namespace conflicts
  \item Use \texttt{r2\_nakagawa()} not \texttt{r2()} — the latter has a bug with certain model structures
  \item Set \texttt{emm\_options(lmer.df = "kenward-roger")} at start of \texttt{emmeans} session
  \item For back-transforming log RT: use \texttt{tran = "log"} in \texttt{emmeans} or exponentiate manually
  \item KR df hang with complex random effects — Satterthwaite is acceptable and standard
  \item Nonword RT: \texttt{emmeans} crashed until \texttt{emm\_options(lmer.df = "satterthwaite")} was set globally
\end{itemize}

% -----------------------------------------------------------------------
\section{5.\ Key Results — Summary}

\subsection{RT Results}

\textbf{Words RT:}
\begin{itemize}[leftmargin=1.5em]
  \item Base Frequency: $\beta = -0.017$, $t(100.97) = -3.66$, $p < .001$, $\eta^2_p = .12$
  \item Family Size: $\beta = -0.013$, $t(93.59) = -2.64$, $p = .010$, $\eta^2_p = .07$
  \item OS: marginal trend, $\beta = -0.035$, $t(63.84) = -1.85$, $p = .069$ — NOT robust to random effects changes
  \item LP: null, $p = .35$
  \item $R^2_\text{marginal} = .033$,\quad $R^2_\text{conditional} = .417$
\end{itemize}

\textbf{Nonwords RT:}
\begin{itemize}[leftmargin=1.5em]
  \item Complexity: $\beta = 0.049$, $t(62.68) = 9.13$, $p < .001$, $\eta^2_p = .57$;
        $M_\text{complex} = 718.9$ ms [695.2, 743.4], $M_\text{simple} = 684.6$ ms [661.5, 708.5]; difference $\approx 34$ ms
  \item Base Frequency: inhibitory effect, $\beta = 0.015$, $t(98.21) = 3.05$, $p = .003$, $\eta^2_p = .09$
  \item Family Size: ns, $p = .22$
  \item OS: significant, $\beta = -0.041$, $t(62.21) = -2.23$, $p = .030$, $\eta^2_p = .07$
  \item LP: null, $p = .78$
  \item $R^2_\text{marginal} = .050$,\quad $R^2_\text{conditional} = .494$
\end{itemize}

\subsection{N250 Results}

\textbf{Words N250 — Base Frequency (primary, tests double dissociation):}
\begin{itemize}[leftmargin=1.5em]
  \item Base Frequency main effect: $\beta = -0.318$, $t(1618) = -3.18$, $p = .002$; High more negative than Low
  \item \textbf{Base Frequency $\times$ OS: $\beta = -0.455$, $t(1618) = -4.06$, $p < .001$} \textit{[Key finding]}
  \item Follow-ups: BF effect absent at low OS ($p = .374$), significant at mean OS ($\Delta = 0.318$, $p = .002$), large at high OS ($\Delta = 0.773$, $p < .001$)
  \item LP: not interacted; main effect ns ($p = .162$)
  \item $R^2_\text{marginal} = .018$,\quad $R^2_\text{conditional} = .801$
\end{itemize}

\textbf{Words N250 — Family Size (secondary):}
\begin{itemize}[leftmargin=1.5em]
  \item Family Size: $\beta = -0.324$, $t(1618) = -5.54$, $p < .001$;
        $M_\text{large} = -0.57\,\mu$V [$-1.03$, $-0.11$], $M_\text{small} = -0.90\,\mu$V [$-1.36$, $-0.44$]
  \item Family Size $\times$ OS: marginal trend, $\beta = -0.111$, $t(1618) = -1.70$, $p = .089$
  \item LP: ns
  \item $R^2_\text{marginal} = .019$,\quad $R^2_\text{conditional} = .738$
\end{itemize}

\textbf{Nonwords N250:}
\begin{itemize}[leftmargin=1.5em]
  \item Complexity: $\beta = -0.233$, $t(4856) = -4.63$, $p < .001$;
        $M_\text{complex} = -0.759\,\mu$V, $M_\text{simple} = -0.526\,\mu$V
  \item Family Size (covariate): $\beta = 0.182$, $t(4856) = 3.63$, $p < .001$
  \item \textbf{Complexity $\times$ OS: $\beta = 0.392$, $t(4856) = 6.98$, $p < .001$} \textit{[Key finding]}
  \item \textbf{Complexity $\times$ LP: $\beta = 0.358$, $t(4856) = 8.76$, $p < .001$} \textit{[Key finding]}
  \item Both interactions show same pattern: large complexity cost at low skill $\to$ attenuated $\to$ reversed at high skill:
  \begin{itemize}
    \item OS: $\Delta = 0.642$ (low), $0.250$ (mean), $-0.142$ (high, $p = .053$)
    \item LP: $\Delta = 0.574$ (low), $0.216$ (mean), $-0.141$ (high, $p = .032$)
  \end{itemize}
  \item Theoretical interpretation: N250 reflects processing cost not sensitivity; reversal = real suffix facilitates parsing for skilled readers; both dimensions converge on early processing (qualifies predicted OS-specific effect)
\end{itemize}

\subsection{N400 Results}

\textbf{Words N400 — Base Frequency (primary, tests double dissociation):}
\begin{itemize}[leftmargin=1.5em]
  \item Base Frequency main effect: $\beta = 0.214$, $t(1618) = 2.01$, $p = .044$; High more positive than Low (facilitation)
  \item \textbf{Base Frequency $\times$ LP: $\beta = 0.344$, $t(1618) = 3.98$, $p < .001$} \textit{[Key finding]}
  \item Follow-ups: BF effect absent at low LP ($p = .333$), significant at mean LP ($\Delta = -0.214$, $p = .044$), large at high LP ($\Delta = -0.558$, $p < .001$)
  \item OS: main effect significant $p = .049$ (overall amplitude, no interaction with BF)
  \item $R^2_\text{marginal} = .026$,\quad $R^2_\text{conditional} = .825$
\end{itemize}

\textbf{Words N400 — Family Size (secondary):}
\begin{itemize}[leftmargin=1.5em]
  \item Family Size: $\beta = -0.452$, $t(1618) = -7.07$, $p < .001$;
        $M_\text{large} = 0.518\,\mu$V [$-0.054$, $1.089$], $M_\text{small} = 0.062\,\mu$V [$-0.509$, $0.634$]
  \item Family Size $\times$ LP: ns, $p = .172$
  \item OS: main effect significant, $p = .038$
\end{itemize}

\textbf{Nonwords N400:}
\begin{itemize}[leftmargin=1.5em]
  \item Complexity: $\beta = -0.460$, $t(4854) = -7.48$, $p < .001$
  \item \textbf{Complexity $\times$ Family Size $\times$ LP three-way: $\beta = 0.227$, $t(4854) = 2.27$, $p = .023$} \textit{[Centerpiece finding]}
  \item Complexity $\times$ Family Size interaction by LP level:
  \begin{itemize}
    \item Low LP: $\Delta = -0.667$, $p < .001$ (large, significant)
    \item Mean LP: $\Delta = -0.440$, $p = .0003$ (attenuated, significant)
    \item High LP: $\Delta = -0.213$, $p = .188$ (absent)
  \end{itemize}
  \item Family size effect within complex nonwords disappears at high LP
  \item Complexity $\times$ Base Frequency $\times$ LP three-way: NOT significant (confirms specificity — LP effect is about family richness, not stem familiarity)
  \item Complexity $\times$ LP (averaged over FS): monotonic attenuation (unlike N250 reversal), significant at all LP levels
\end{itemize}

\subsection{The Double Dissociation — Confirmed}

\begin{center}
\begin{tabular}{lcc}
  \toprule
   & \textbf{N250 (early form-based)} & \textbf{N400 (later semantic)} \\
  \midrule
  \textbf{Base Frequency $\times$ OS} & \done\ Significant & OS main effect only \\
  \textbf{Base Frequency $\times$ LP} & LP not interacted  & \done\ Significant \\
  \bottomrule
\end{tabular}
\end{center}

\smallskip
For nonwords: both OS and LP modulate N250 (partial qualification of predicted dissociation); LP selectively drives N400 three-way.

% -----------------------------------------------------------------------
\section{6.\ Results Section — Current Status}

\subsection{Written and Approved Paragraphs}

\textbf{Words section:}
\begin{itemize}[leftmargin=1.5em]
  \item[\done] Words RT — complete, verified against R output
  \item[\done] Words N250 (family size version) — complete
  \item[\done] Words N400 (family size version) — complete
  \item[\done] Words N250 (base frequency version) — drafted in final exchange
  \item[\done] Words N400 (base frequency version) — drafted in final exchange
\end{itemize}

\textbf{Nonwords section:}
\begin{itemize}[leftmargin=1.5em]
  \item[\done] Nonwords RT — complete, verified against R output
  \item[\done] Nonwords N250 — complete with follow-up contrasts
  \item[\done] Nonwords N400 — complete with three-way follow-up contrasts
\end{itemize}

\textbf{Summary section:}
\begin{itemize}[leftmargin=1.5em]
  \item[\done] Results Summary — drafted and approved, split so that the ``three key findings'' paragraph moves to open the Discussion
\end{itemize}

\subsection{Results Section Structure}

\begin{lstlisting}[language={}, style=Rstyle, basicstyle=\ttfamily\small]
\section{Results}
  [Orienting paragraph]
  [Table 1: Model summary]
  [Asymmetric interaction explanation paragraph]

\subsection{Words}
  \paragraph{Words RT}
  \paragraph{Words N250}   % <- base frequency version now primary
  \paragraph{Words N400}   % <- base frequency version now primary

\subsection{Nonwords}
  \paragraph{Nonwords RT}
  \paragraph{Nonwords N250}
  \paragraph{Nonwords N400}

\subsection{Results Summary}
\end{lstlisting}

\subsection{Key Writing Decisions}

\begin{itemize}[leftmargin=1.5em]
  \item Words and nonwords reported separately (not by component)
  \item RT then ERP (N250 then N400) within each stimulus type
  \item No claim that ``interactions were absent in RTs'' since they were not tested — instead state models were specified for efficiency only
  \item ERP measures framed as ``designed to detect'' modulation that RT models were not designed to detect
  \item N250 direction for base frequency: high frequency = MORE negative (opposite to family size effect) — to be explained in Discussion
  \item Conditional $R^2$ values for ERP models are very high (0.738--0.825) due to large participant and electrode variance
\end{itemize}

% -----------------------------------------------------------------------
\section{7.\ Outstanding Issues to Resolve}

\subsection{Critical (resolve before opening new chat)}

\begin{enumerate}[leftmargin=1.5em]
  \item \textbf{Nonword N250 covariate} — \textbf{RESOLVED:} Final model uses \texttt{family\_size} as covariate.

  \item \textbf{Integration of base frequency and family size analyses for words:} The words section now has two sets of ERP analyses — base frequency (primary, tests double dissociation) and family size (secondary). Need to decide:
  \begin{itemize}
    \item Report base frequency as primary, family size as secondary/supplementary?
    \item Or report both in full within the words section?
    \item \textit{Recommended:} base frequency primary in main text, family size results noted briefly or in supplementary.
  \end{itemize}

  \item \textbf{N250 direction for base frequency:} High frequency bases $\to$ MORE negative N250 (opposite of family size effect). This dissociation between BF and FS at N250 needs a sentence in Discussion. Likely explanation: BF = strength of form-level activation (more effort); FS = efficiency of pattern matching (less effort).
\end{enumerate}

\subsection{Important (address in new chat)}

\begin{enumerate}[leftmargin=1.5em, start=4]
  \item \textbf{Discussion — not yet started.} Structure agreed:
  \begin{itemize}
    \item Opening: three key findings paragraph (moved from Results Summary)
    \item Section 1: N250 reversal at high skill (most unexpected finding)
    \item Section 2: Behavioral/ERP dissociation
    \item Section 3: N400 three-way and semantic integration (connect to theoretical frameworks)
    \item Closing: theoretical implications and limitations
  \end{itemize}

  \item \textbf{Updated model summary table:} Table 1 needs updating to reflect base frequency as primary predictor for words ERP analyses (currently shows family size).

  \item \textbf{Results orienting paragraph:} Needs minor update to reflect that base frequency analyses are primary for words.

  \item \textbf{Nonword N400 base frequency analyses:} Null three-way (Complexity $\times$ BF $\times$ LP) should be reported as a cross-check confirming specificity of the family size finding — currently mentioned but not formally integrated into the results write-up.
\end{enumerate}

% -----------------------------------------------------------------------
\section{8.\ Discussion — Planned Structure (not yet written)}

\subsection{Opening Paragraph (moved from Results Summary)}

The three key findings that structure the Discussion:
\begin{enumerate}[leftmargin=1.5em]
  \item Individual differences dissociate from behavioral responses but are visible in neural dynamics
  \item LP selectively drives the N400 three-way (closest support for predicted dissociation)
  \item N250 shows reversal at high skill; N400 shows attenuation without reversal — qualitatively different mechanisms
\end{enumerate}

\subsection{Section 1: Early Morphological Parsing and the N250 Reversal}
\begin{itemize}[leftmargin=1.5em]
  \item Processing cost vs.\ morphological sensitivity interpretation
  \item Why the reversal indicates real suffix facilitates parsing for skilled readers
  \item Why both OS and LP converge on early processing — broad lexical experience, not just orthographic skill
  \item Marginal Family Size $\times$ OS for words — consistency with nonword pattern
  \item Connect to morpho-orthographic and discriminative learning accounts
\end{itemize}

\subsection{Section 2: The Behavioral/ERP Dissociation}
\begin{itemize}[leftmargin=1.5em]
  \item Individual differences invisible in RTs but visible in ERPs
  \item Implications for measuring individual differences in morphological processing
  \item Lexical Quality Hypothesis — precision of representation vs.\ speed of response
  \item Why OS predicts nonword RT but not morphological modulation in ERPs
\end{itemize}

\subsection{Section 3: N400 Three-Way and Semantic Integration}
\begin{itemize}[leftmargin=1.5em]
  \item What Complexity $\times$ Family Size $\times$ LP means
  \item Attenuation without reversal — contrast with N250
  \item LP selective at N400; OS not — closest to predicted dissociation
  \item Connect to NDL/LDL: LP indexes depth of form-to-meaning mapping
  \item Connect to multi-stage accounts
  \item Null Family Size $\times$ LP for words — why LP does not modulate FS for familiar words
  \item Base Frequency $\times$ LP at N400 for words — LP enhances semantic integration benefit of high frequency stems
\end{itemize}

\subsection{Closing: Theoretical Implications and Limitations}
\begin{itemize}[leftmargin=1.5em]
  \item Partial dissociation — what it means
  \item Limitations: unmatched words, asymptotic CIs, generalisability
  \item Future directions
\end{itemize}

% -----------------------------------------------------------------------
\section{9.\ Key Terminology and Variable Names}

\begin{center}
\begin{tabular}{p{4.2cm} p{5.8cm} p{4.5cm}}
  \toprule
  \textbf{Concept} & \textbf{Variable name in R} & \textbf{Notes} \\
  \midrule
  Language Proficiency     & \texttt{LP}, \texttt{Dim.1}      & PCA dimension 1, already centred \\[2pt]
  Orthographic Sensitivity & \texttt{OS}, \texttt{Dim.2}      & PCA dimension 2, already centred \\[2pt]
  Base/Stem Frequency      & \texttt{zLogBF} (RT), \texttt{base\_frequency} (ERP) & Continuous for RT, binned for ERP \\[2pt]
  Family Size              & \texttt{zLogFS} (RT), \texttt{family\_size} (ERP)    & Continuous for RT, binned for ERP \\[2pt]
  Complexity               & \texttt{complexity}              & Factor: Simple vs.\ Complex \\[2pt]
  Participant ID           & \texttt{SubjID}                  & \\[2pt]
  Item ID                  & \texttt{Item} (words), \texttt{ItemID} (nonwords) & Note different names \\[2pt]
  ERP amplitude            & \texttt{value}                   & \\[2pt]
  Electrode nested in participant & \texttt{SubjID:chlabel}   & \\[2pt]
  Log RT                   & \texttt{logRT}                   & \texttt{log(response\_time)} \\
  \bottomrule
\end{tabular}
\end{center}

% -----------------------------------------------------------------------
\section{10.\ Citations Needing Completion}

The following citations appear as \texttt{[CITATION]} placeholders throughout the manuscript:

\begin{itemize}[leftmargin=1.5em]
  \item Baayen \& Schreuder race model — original dual-route paper
  \item Rastle, Davis et al.\ — early automatic decomposition / masked priming papers
  \item Grainger et al.\ — morpho-orthographic account
  \item Baayen et al.\ NDL paper — 2011 \textit{Psychological Review}
  \item Baayen et al.\ LDL paper — \textasciitilde2019 \textit{Complexity} journal
  \item Family size effects — original Schreuder \& Baayen paper
  \item Stem frequency operationalisation — relevant citations
  \item Semantic transparency and family size — opacity studies
  \item N250 and N400 components — relevant ERP citations
  \item Andrews et al.\ (2013) — individual differences in morphological priming
  \item Perfetti — Lexical Quality Hypothesis (2002, 2007, 2017)
  \item Andrews and colleagues — lexical precision papers
  \item Marslen-Wilson, Tyler, Devlin — neuroimaging morphology papers
\end{itemize}

\bigskip
\hrule
\smallskip
{\small \textit{Log created at end of extended working conversation. Open new chat in this project and paste this document as the first message to establish context before posting new R output.}}

\end{document}
